\uuid{3Aya}
\titre{Limite en un point}
\auteur{bodin ; hanani ; mzari}
\organisation{exo7}
\licence{CC BY-NC-SA 4.0}
\datecreate{2020-01-20}
\niveau{L1}
\module{Analyse}
\chapitre{Continuité, limite et étude de fonctions réelles}
\sousChapitre{Limite de fonctions}
\difficulte{3}
\theme{Limites des fonctions réelles, limite en un point}
\exo7qcmid{lille-1.limites.026}

\contenu{
\begin{question}
\qclassification{123.03}
\qtags{motcle=limite en un point}
Soit $f(x)= \frac{\sqrt{x+1}-1}{\sqrt[3]{x+1}-1}$. Quelles sont les assertions vraies ?
\begin{answers}

    \bad{$\lim_{x\to 0} f(x)=0$}

    \good{$\lim_{x\to 0} f(x)=\frac{3}{2}$}

    \bad{$f$ n'admet pas de limite en $0$.}

    \bad{$\lim_{x\to 0} f(x)=+\infty$}
\end{answers}
\begin{explanations}
On pourra multiplier $f$ par  $\sqrt{x+1}+1$ et $(\sqrt[3]{x+1})^2+\sqrt[3]{x+1}+1$ les expressions conjuguées de $\sqrt{x+1}-1$ et de $\sqrt[3]{x+1}-1$ respectivement. On obtient :
$f(x)=\frac{(\sqrt[3]{x+1})^2+\sqrt[3]{x+1}+1}{\sqrt{x+1}+1}$.
\end{explanations}
\end{question}
}
